\section{Methodology}
\label{sec:methodology}

\subsection{Overview}

Given a source code function $f$, VulGCL produces a binary label
$\hat{y} \in \{0,1\}$ (non-vulnerable / vulnerable) by fusing three
independently computed representations:

\begin{equation}
\hat{y} = \text{Classifier}\bigl(
  \mathbf{h}_G \,\|\, \mathbf{h}_I \,\|\, \mathbf{h}_L
\bigr)
\end{equation}

where $\mathbf{h}_G$, $\mathbf{h}_I$, $\mathbf{h}_L$ are the graph, image, and
LLM branch embeddings respectively, and $\|$ denotes concatenation.

\subsection{PDG Extraction}
\label{sec:pdg}

We parse each function using Joern~\cite{joern} to extract its Program
Dependency Graph $\mathcal{G} = (\mathcal{V}, \mathcal{E})$, where each node
$v \in \mathcal{V}$ represents a code statement and each edge
$e \in \mathcal{E}$ represents a data or control dependency between statements.
Node features are initialized as the sentence embedding of the corresponding
code statement using CodeBERT tokenization.

\subsection{Graph Branch ($\mathbf{h}_G$)}
\label{sec:graph_branch}

% TODO: describe GNN architecture (GCN / GAT / GIN — decide after ablation)
The PDG $\mathcal{G}$ is passed through $L$ layers of a Graph Attention Network
(GAT):

\begin{equation}
\mathbf{h}_G = \text{Readout}\Bigl(\text{GAT}^{(L)}(\mathcal{G})\Bigr)
\end{equation}

Global mean pooling is used as the readout function.

\subsection{Image Branch ($\mathbf{h}_I$)}
\label{sec:image_branch}

Following VulCNN~\cite{vulcnn}, we encode each node $v$ in the PDG as a
sentence vector $\mathbf{s}_v \in \mathbb{R}^d$ (sent2vec), then weight each
row by three centrality measures computed on the PDG:

\begin{itemize}
  \item Channel 1: Degree centrality $\kappa_{\text{deg}}(v)$
  \item Channel 2: Katz centrality $\kappa_{\text{katz}}(v)$
  \item Channel 3: Closeness centrality $\kappa_{\text{clo}}(v)$
\end{itemize}

The result is a three-channel image $\mathbf{I} \in \mathbb{R}^{N \times d \times 3}$
passed to a CNN:

\begin{equation}
\mathbf{h}_I = \text{CNN}(\mathbf{I})
\end{equation}

\subsection{LLM Branch ($\mathbf{h}_L$)}
\label{sec:llm_branch}

The raw token sequence of the function is fed into CodeBERT~\cite{codebert}.
We take the \texttt{[CLS]} representation as the function-level semantic
embedding:

\begin{equation}
\mathbf{h}_L = \text{CodeBERT}(f)[\texttt{CLS}]
\end{equation}

CodeBERT weights are fine-tuned end-to-end during training.

\subsection{Fusion and Classification}

The three embeddings are concatenated and passed through a two-layer MLP
classifier with dropout:

\begin{equation}
\hat{y} = \sigma\Bigl(W_2 \cdot \text{ReLU}(W_1 \cdot [\mathbf{h}_G \| \mathbf{h}_I \| \mathbf{h}_L] + b_1) + b_2\Bigr)
\end{equation}

We train with binary cross-entropy loss.

\subsection{VulGCL-Lite: Edge Deployment}
\label{sec:deployment}

To enable inference on a Raspberry Pi 4B (ARM Cortex-A72, 8\,GB RAM), we apply
post-training quantization: weights and activations are converted from
float32 to int8 using ONNX Runtime. The LLM branch is further compressed via
knowledge distillation from CodeBERT to a 2-layer LSTM student encoder.

% TODO: Add architecture diagram (Figure 1)

\subsection{Dataset and Labeling}
\label{sec:data}

\textbf{Standard benchmarks.} We evaluate on Devign~\cite{devign}
(27,318 functions, C) and BigVul~\cite{bigvul} (188,636 functions, C/C++),
using the standard train/val/test splits.

\textbf{Stack Overflow supplementary corpus.} From our 620K-post SO dataset
(2015--2025), we extract Python and Java code snippets from accepted answers.
Snippets are labeled using Semgrep static analysis. Raw vote scores are
\emph{not} used as labels due to Chen et al.'s~\cite{chen2019reliable} finding
that insecure SO answers average higher scores than secure ones. Instead, we use
a composite signal: static analysis label + comment-flag pattern (e.g.,
\emph{``this is insecure''}) + accepted-answer status.
